% =====================================================================
%  CHAPTER 5: 社会环境因素与健康
% =====================================================================
\chapter{社会环境因素与健康}
\setcounter{chapter}{5}
\chapterintro{
掌握人口结构与健康、人口流动与健康的关系（\textbf{教学难点}）；掌握社会支持的概念及类型；熟悉食品安全概念及防治、家庭对健康的影响；理解社会阶层与健康的关系；了解人口数量、营养、体力活动、城市化等因素的健康影响。
}

% ----- 5.1 人口与健康 -----
\section{人口与健康}

\subsection{人口数量与健康}

\begin{defbox}
\textbf{人口数量与健康的关系：}
\begin{itemize}[leftmargin=*]
    \item 人口过多→资源（粮食、水、能源）紧张→营养不良、环境污染、居住拥挤→传染病易于传播
    \item 人口过少→劳动力不足→经济发展受限→卫生投入不足
    \item \imp{人口密度}过高增加传染病传播风险和心理压力
    \item 世界人口持续增长对全球卫生系统构成持续性压力
\end{itemize}
\end{defbox}

\subsection{人口结构与健康（教学难点*）}

\begin{keybox}
\imp{（教学难点*）}\textbf{人口年龄结构与健康：}
\begin{enumerate}[leftmargin=*]
    \item \imp{年龄结构类型}：
    \begin{itemize}
        \item 年轻型（增长型）——0-14岁人口占比$\geq$40\%，65岁及以上$\leq$4\%
        \item 成年型（稳定型）——0-14岁占30-40\%，65岁及以上占4-7\%
        \item 老年型（缩减型）——0-14岁$\leq$30\%，65岁及以上$\geq$7\%
    \end{itemize}
    \item \imp{人口老龄化}：60岁以上人口占总人口比例$\geq$10\%或65岁以上$\geq$7\%即为老龄化社会

    \item \textbf{老龄化的健康影响：}
    \begin{itemize}
        \item 慢性非传染性疾病负担加重（心脑血管病、肿瘤、糖尿病等）
        \item 老年综合征（跌倒、失能、痴呆、抑郁等）增加
        \item 长期照护需求急剧增长
        \item 社会保障和医疗卫生系统压力增大
    \end{itemize}

    \item \imp{抚养系数}：
    \begin{itemize}
        \item 总抚养系数 = [(14岁及以下人口数 + 65岁及以上人口数) / 15-64岁人口数] $\times$ 100\%
        \item 少儿抚养系数 = (14岁及以下人口数 / 15-64岁人口数) $\times$ 100\%
        \item 老年抚养系数 = (65岁及以上人口数 / 15-64岁人口数) $\times$ 100\%
    \end{itemize}

    \item \imp{性别结构}：出生性别比 = (某年活产男婴数 / 同年活产女婴数) $\times$ 100，正常范围为103-107。性别比失调会带来一系列社会健康和稳定问题
\end{enumerate}

\textbf{人口金字塔的三种类型：}
\begin{itemize}[leftmargin=*]
    \item \imp{增长型（年轻型）}——塔基底部宽、顶部尖，年轻人比重大，预示着人口将持续快速增长
    \item \imp{稳定型（成年型）}——塔底和塔身宽度基本接近，塔尖逐渐收缩，各年龄段人口比例相对均衡
    \item \imp{缩减型（老年型）}——塔底窄、塔身宽，年轻人口少、中老年人口比例大，预示着人口将逐渐减少
\end{itemize}
\end{keybox}

\subsection{人口素质与健康}

\begin{defbox}
\textbf{人口素质}包括三个方面：
\begin{enumerate}[leftmargin=*]
    \item \imp{思想道德素质}
    \item \imp{科学文化素质}——受教育水平直接影响健康知识获取和健康行为养成
    \item \imp{身体素质}——包括体格发育、生理机能、智力水平等，与遗传、营养、环境有关
\end{enumerate}

\textbf{出生缺陷}是影响人口素质的重要因素——我国出生缺陷发生率约为5.6\%，是婴儿死亡和儿童残疾的首要原因，带来沉重的社会疾病负担。
\end{defbox}

\subsection{人口流动与健康（教学难点*）}

\begin{keybox}
\imp{（教学难点*）}\textbf{人口流动对健康的影响：}

\textbf{流动人口面临的主要健康问题：}
\begin{enumerate}[leftmargin=*]
    \item \imp{传染病传播风险}：流动人口密集的生活和工作条件有利于传染病的传播（结核病、性传播疾病、艾滋病等）
    \item \imp{职业健康危害}：多从事高风险职业（建筑、采矿、制造业等），职业伤害和职业病发病率高
    \item \imp{妇幼健康问题}：流动妇女产前检查率低、住院分娩率低，新生儿和婴儿死亡率较高
    \item \imp{心理健康问题}：社会隔离、歧视、工作压力、家庭分离等导致心理问题
    \item \imp{卫生服务可及性差}：医疗保险关系转移接续困难，经济负担重，就医不便
    \item \imp{儿童健康与发展}：留守儿童和流动儿童的健康、营养、心理和教育问题突出
    \item \imp{"输入性"传染病风险}：某些地区性传染病可能随人口流动传播到新地区
\end{enumerate}
\end{keybox}

% ----- 5.2 营养、食品安全与健康 -----
\section{营养、食品安全与健康}

\subsection{营养与健康}

\begin{enumerate}[leftmargin=*]
    \item \imp{营养不良}——包括营养不足（消瘦、发育迟缓、微量营养素缺乏）和营养过剩（超重、肥胖）。我国面临营养不良的\keyword{双重负担}
    \item \imp{营养过剩}——导致肥胖、2型糖尿病、心血管疾病、某些癌症的风险增加
    \item \imp{膳食结构不合理}——高盐、高脂、低蔬菜水果摄入增加慢性病风险
\end{enumerate}

\subsection{食品安全}

\begin{defbox}
\textbf{食品安全（Food Safety）}：指食品无毒、无害，符合应当有的营养要求，对人体健康不造成任何急性、亚急性或者慢性危害。
\end{defbox}

\begin{keybox}
\textbf{食品安全事件的危害：}
\begin{enumerate}[leftmargin=*]
    \item 急性食源性疾病——微生物污染（细菌、病毒）导致食物中毒
    \item 慢性健康损害——农药残留、重金属、食品添加剂等长期低剂量暴露
    \item 社会恐慌和社会心理影响
    \item 经济损失——影响食品产业和国际贸易
\end{enumerate}

\textbf{食品安全的社会防治：}
\begin{enumerate}[leftmargin=*]
    \item 完善食品安全法律法规和标准体系（《食品安全法》）
    \item 建立食品安全全程监管体系——从农田到餐桌
    \item 加强食品安全风险监测与评估
    \item 提高食品安全检测技术和能力
    \item 加强食品安全教育和信息沟通
    \item 鼓励社会监督和参与
\end{enumerate}
\end{keybox}

% ----- 5.3 体力活动与职业压力 -----
\section{体力活动、职业压力与健康}

\subsection{体力活动与健康}

\begin{enumerate}[leftmargin=*]
    \item \imp{缺乏体力活动}已成为全球第四大死亡危险因素
    \item 适量体力活动可降低心血管疾病、糖尿病、某些癌症、抑郁症的风险
    \item WHO建议成人每周至少150分钟中等强度有氧运动
    \item 久坐生活方式（Sedentary Lifestyle）是独立的健康危险因素
\end{enumerate}

\subsection{职业压力与健康}

\begin{defbox}
\textbf{职业压力（Occupational Stress）}：当工作要求与工作者的能力、资源或需求不匹配时产生的有害的生理和情绪反应。
\end{defbox}

\textbf{职业压力的健康影响：}心血管疾病、心理障碍（抑郁、焦虑、职业倦怠）、免疫功能下降、肌肉骨骼疾患增加。

% ----- 5.4 城市化与健康 -----
\section{城市化与健康}

\begin{defbox}
\textbf{城市化（Urbanization）}：由农业为主的传统乡村社会向以工业和服务业为主的现代城市社会逐渐转变的历史过程。
\end{defbox}

\begin{keybox}
\textbf{城市化对健康的影响：}

\textbf{积极影响：}
\begin{enumerate}[leftmargin=*]
    \item 卫生服务可及性更好
    \item 教育机会更多→健康素养更高
    \item 就业机会和收入水平较高
    \item 基础设施（供水、排污、交通）改善
\end{enumerate}

\textbf{消极影响（"城市病"）：}
\begin{enumerate}[leftmargin=*]
    \item 环境污染（空气、噪声、光污染）
    \item 生活节奏快、竞争压力大→心理问题
    \item 慢性病高发（饮食西方化、久坐生活方式）
    \item 交通拥堵和交通事故
    \item 社会关系疏离→社会支持减弱
    \item 传染病在密集人口中的传播风险
    \item "城市贫民窟"问题——居住条件恶劣、卫生设施缺乏
\end{enumerate}
\end{keybox}

% ----- 5.5 社会网络与社会支持 -----
\section{社会网络与社会支持}

\subsection{社会网络}

\begin{defbox}
\textbf{社会网络（Social Network）}：个人与其他人之间所有社会联系所构成的网络系统。社会网络以个体为中心，围绕个体的社会关系形成。
\end{defbox}

\begin{itemize}[leftmargin=*]
    \item \imp{社会网络的功能}：①提供社会支持；②传递信息和价值观；③影响健康行为；④提供社会参与机会
    \item 社会网络的密度、规模、强度和互惠性影响其对健康的保护作用
\end{itemize}

\subsection{人际关系与社会支持}

\begin{defbox}
\textbf{人际关系（Interpersonal Relationship）}：人们在社会交往过程中形成的各种关系的总称，包括家庭关系、朋友关系、同事关系、邻里关系等。
\end{defbox}

\begin{keybox}
\textbf{社会支持（Social Support）}：个人从社会网络中所获得的物质和精神的帮助与支持。

\textbf{社会支持的类型：}
\begin{enumerate}[leftmargin=*]
    \item \imp{情感支持（Emotional Support）}：提供共情、关心、爱护与信任，帮助个体应对心理压力
    \item \imp{物质支持（Instrumental Support）}：提供实际的物质帮助和服务（金钱、物品、人力帮助等）
    \item \imp{信息支持（Informational Support）}：提供有助于解决问题的建议、指导和信息
    \item \imp{评价支持（Appraisal Support）}：提供反馈、肯定和社会比较，帮助个体进行自我评价
\end{enumerate}

\textbf{社会支持对健康的影响机制：}
\begin{enumerate}[leftmargin=*]
    \item \imp{直接效应（主效应模型）}：良好的社会支持直接促进健康，无论是否处于应激状态
    \item \imp{缓冲效应（压力缓冲模型）}：社会支持缓冲压力事件对健康的负面冲击
    \item 社会支持可改善免疫功能、降低心血管反应性、促进健康行为
\end{enumerate}
\end{keybox}

% ----- 5.6 家庭与健康 -----
\section{家庭与健康}

\subsection{家庭的类型与功能}

\begin{defbox}
\textbf{家庭}是以婚姻和血缘关系为基础建立起来的社会生活的基本单位。

\textbf{家庭类型：}
\begin{itemize}[leftmargin=*]
    \item 核心家庭——由父母和未婚子女组成
    \item 扩展家庭——由两代或两代以上的夫妇及子女组成
    \item 单亲家庭
    \item 其他特殊家庭（重组家庭、空巢家庭等）
\end{itemize}

\textbf{家庭功能：}
\begin{enumerate}[leftmargin=*]
    \item 生育功能
    \item 抚养和赡养功能
    \item 经济功能
    \item 情感支持功能
    \item 社会化功能——价值观、行为规范的内化
    \item 健康照顾功能——为家庭成员提供工具性和情感性健康支持
\end{enumerate}

\textbf{家庭功能评估——Family APGAR问卷：}
\begin{itemize}[leftmargin=*]
    \item A（Adaptation，适应度）——家庭遭遇危机时，利用内外资源解决问题的能力
    \item P（Partnership，合作度）——家庭成员分担责任和共同作出决定的程度
    \item G（Growth，成长度）——家庭成员通过互相支持所达到的身心成熟和自我实现程度
    \item A（Affection，情感度）——家庭成员间相互关爱、表达情感的程度
    \item R（Resolve，亲密度）——家庭成员共享时间、空间和金钱的程度
\end{itemize}
\end{defbox}

\subsection{家庭对健康的影响}

\begin{keybox}
\textbf{家庭对健康的积极影响：}
\begin{enumerate}[leftmargin=*]
    \item 提供情感支持——减少心理问题和精神疾病
    \item 影响健康行为的形成——饮食、运动、卫生习惯等
    \item 在疾病时提供照护
    \item 提供经济和物质保障
\end{enumerate}

\textbf{家庭对健康的消极影响：}
\begin{enumerate}[leftmargin=*]
    \item 家庭暴力——直接躯体伤害和心理创伤
    \item 家庭功能失调——增加成员患病风险
    \item 家庭压力（离婚、丧偶、经济困难等）——增加心理和躯体疾病风险
    \item 家庭中传染病传播（如结核病、肝炎等）
\end{enumerate}
\end{keybox}

% ----- 5.7 退休、失业与健康 -----
\section{退休、失业与健康}

\begin{enumerate}[leftmargin=*]
    \item \imp{失业与健康}：失业导致收入下降、社会地位丧失、心理压力增加，与死亡率升高、心理健康问题、自杀风险增加密切相关
    \item \imp{退休与健康}：退休是人生重要的角色转换，适应良好者健康不受影响甚至改善，适应不良者可能出现"退休综合征"——抑郁、失落感、躯体不适等
\end{enumerate}

% ----- 5.8 社会阶层与健康 -----
\section{社会阶层与健康}

\begin{defbox}
\textbf{社会阶层（Social Class/Social Stratum）}：由于经济、政治、社会等多种原因形成的，在社会层次结构中处于不同地位的社会群体。划分社会阶层的三大指标：\imp{经济地位/收入}（最主要）、\imp{职业}、\imp{教育水平}。

\textbf{中国的十大社会阶层（中国社会科学院划分）：}
\begin{enumerate}[leftmargin=*]
    \item 国家与社会管理者阶层
    \item 经理人员阶层
    \item 私营企业主阶层
    \item 专业技术人员阶层
    \item 办事人员阶层
    \item 个体工商户阶层
    \item 商业服务业员工阶层
    \item 产业工人阶层
    \item 农业劳动者阶层
    \item 城乡无业、失业、半失业者阶层
\end{enumerate}
\end{defbox}

\begin{keybox}
\textbf{社会阶层对健康的影响（健康梯度现象）：}
\begin{enumerate}[leftmargin=*]
    \item \imp{健康梯度（Health Gradient）}：社会阶层每上升一个层级，健康状况和期望寿命均有改善，形成从高到低的梯度关系，而非仅"两极"差异
    \item 低社会阶层人群具有更高的死亡率、发病率和更低的期望寿命
    \item \textbf{影响机制：}
    \begin{itemize}
        \item 物质路径：收入差异→生活条件、营养、医疗可及性差异
        \item 心理路径：低社会地位→相对剥夺感、失控感→慢性应激
        \item 行为路径：健康素养差异→健康行为差异
        \item 生命周期累积：生命历程中社会劣势的累积效应
    \end{itemize}
    \item \textbf{Wilkinson假说}：在富裕国家，社会不平等程度（而非平均收入水平）是决定人群健康水平的关键因素
\end{enumerate}
\end{keybox}

% ----- 复习题 -----
\section{复习题}

\subsection*{一、选择题}
\begin{enumerate}[leftmargin=*, itemsep=3pt]
    \item 衡量人口老龄化最常用的指标是：
    \begin{enumerate}
        \item[(A)] 60岁以上人口$\geq$10\%或65岁以上$\geq$7\%
        \item[(B)] 60岁以上人口$\geq$15\%或65岁以上$\geq$10\%
        \item[(C)] 50岁以上人口$\geq$20\%
        \item[(D)] 70岁以上人口$\geq$5\%
    \end{enumerate}
    \item 社会支持的类型不包括：
    \begin{enumerate}
        \item[(A)] 情感支持
        \item[(B)] 物质支持
        \item[(C)] 政治支持
        \item[(D)] 信息支持
    \end{enumerate}
    \item 以下属于"健康梯度"现象的是：
    \begin{enumerate}
        \item[(A)] 只有最贫困和最富有人群的健康存在差异
        \item[(B)] 社会阶层每上升一级，健康状况均有改善
        \item[(C)] 社会阶层与健康完全无关
        \item[(D)] 中产阶级的健康状况最差
    \end{enumerate}
    \item 正常出生性别比的范围是：
    \begin{enumerate}
        \item[(A)] 95-100
        \item[(B)] 103-107
        \item[(C)] 110-115
        \item[(D)] 100-103
    \end{enumerate}
\end{enumerate}

\subsection*{二、名词解释}
\begin{enumerate}[leftmargin=*, itemsep=3pt]
    \item \textbf{社会支持（Social Support）}
    \item \textbf{健康梯度（Health Gradient）}
    \item \textbf{城市化（Urbanization）}
    \item \textbf{人口老龄化}
\end{enumerate}

\subsection*{三、简答题}
\begin{enumerate}[leftmargin=*, itemsep=3pt]
    \item 试述人口流动对健康的影响。
    \item 简述社会支持的类型及其对健康的影响机制。
    \item 简述社会阶层对健康的影响及可能机制。
    \item 简述食品安全的概念及其社会防治措施。
    \item 简述城市化对健康的正反两面影响。
\end{enumerate}

\subsection*{参考答案要点}

\subsubsection*{一、选择题}
1. (A)；2. (C)；3. (B)；4. (B)

\subsubsection*{二、名词解释}
\begin{enumerate}[leftmargin=*]
    \item \textbf{社会支持}：个人从社会网络中获得物质和精神帮助与支持的总和，包括情感、物质、信息和评价支持四种类型。
    \item \textbf{健康梯度}：社会阶层与健康水平之间存在的梯度关系，社会阶层每上升一个层级，健康状况均有改善。
    \item \textbf{城市化}：由农业为主的乡村社会向以工业和服务业为主的现代城市社会转变的过程。
    \item \textbf{人口老龄化}：60岁以上人口占总人口比例$\geq$10\%或65岁以上$\geq$7\%的社会现象。
\end{enumerate}

\subsubsection*{三、简答题}
\begin{enumerate}[leftmargin=*]
    \item 七大健康问题：传染病传播风险、职业健康危害、妇幼健康问题、心理健康问题、卫生服务可及性差、儿童健康与发展问题、"输入性"传染病风险。
    \item 四种类型：情感支持（共情和关爱）、物质支持（实际帮助）、信息支持（建议指导）、评价支持（反馈和比较）。影响机制：直接效应（主效应模型）和缓冲效应（压力缓冲模型）。
    \item 健康梯度现象：社会阶层每上升一级，健康均有改善。机制：物质路径（收入和生活条件）、心理路径（相对剥夺感和慢性应激）、行为路径（健康素养和行为差异）、生命周期累积效应。
    \item \textbf{食品安全}：食品无毒无害、符合营养要求、不造成健康危害。防治措施：完善法规和标准、全程监管体系、风险监测评估、提高检测能力、加强教育沟通、鼓励社会监督。
    \item 积极：卫生服务可及性好、教育机会多、收入高、基础设施改善。消极：环境污染、心理压力大、慢性病高发、交通事故、社会关系疏离、"城市贫民窟"问题。
\end{enumerate}

\newpage
